\newglossaryentry{heig-vd}{
    name=HEIG-VD,
    description={Haute École d'Ingénierie et de Gestion du canton de Vaud}
}
\newglossaryentry{kubernetes}{
    name=Kubernetes,
    description={Système d'orchestration de conteneurs open-source permettant d'automatiser
    le déploiement, la mise à l'échelle et la gestion d'applications conteneurisées}
}
\newglossaryentry{ray}{
    name=Ray,
    description={Framework Python de calcul distribué, optimisé pour les workloads
    GPU et les pipelines d'intelligence artificielle}
}
\newglossaryentry{zarr}{
    name=ZARR,
    description={Format de données N-dimensionnel orienté chunking,
    permettant la lecture partielle (par tuiles) de grandes images sans les charger entièrement en mémoire}
}
\newglossaryentry{sam3}{
    name=SAM3,
    description={Segment Anything Model version 3 --- modèle de segmentation d'images
    développé par Meta AI, capable de segmenter n'importe quel objet dans une image}
}
\newglossaryentry{minio}{
    name=MinIO,
    description={Serveur de stockage objet haute performance compatible avec le protocole S3 d'Amazon}
}
\newglossaryentry{postgresql}{
    name=PostgreSQL,
    description={Système de gestion de base de données relationnelle open-source}
}
\newglossaryentry{postgis}{
    name=PostGIS,
    description={Extension de PostgreSQL ajoutant le support des types et fonctions géospatiales,
    permettant de stocker et requêter des géométries (points, polygones, etc.)}
}
\newglossaryentry{labelstudio}{
    name={Label Studio},
    description={Plateforme open-source d'annotation de données pour les projets de machine learning,
    supportant images, texte, audio et vidéo}
}
\newglossaryentry{pod}{
    name=Pod,
    description={Unité de déploiement atomique dans Kubernetes, regroupant un ou plusieurs conteneurs}
}
